\documentclass{chsmc}
\chsmcsetup{
  year = 2021,
  part = 1,
  type = A,
  show-answers = false,
}
\begin{document}

\section{填空题：本大题共 8 小题，每小题 8 分，共 64 分.}

\begin{enumerate}
  \item 等差数列 $\{a_n\}$ 满足 $a_{2021} = a_{20} + a_{21} = 1$，
    则 $a_1$ 的值为 \blankline
  \item 设集合 $A = \{1, 2, m\}$，其中 $m$ 为实数. 令
    $B = \{a^2 \mid a \in A\}$，$C = A \cup B$. 若 $C$ 的所有元素之和为 $6$，
    则 $C$ 的所有元素之积为 \blankline
  \item 设函数 $f(x)$ 满足：对任意非零实数 $x$，均有
    $f(x) = f(1)\cdot x + \dfrac{f(2)}{x} - 1$，
    则 $f(x)$ 在 $(0,+\infty)$ 上的最小值为 \blankline
  \item 设函数 $f(x) = \cos x + \log_2 x$ $(x > 0)$，若正实数 $a$ 满足
    $f(a) = f(2a)$，则 $f(2a) - f(4a)$ 的值为 \blankline
  \item 在 $\triangle ABC$ 中，$AB = 1$，$AC = 2$，
    $B - C = \dfrac{2\pi}{3}$，则 $\triangle ABC$ 的面积为 \blankline
  \item 在平面直角坐标系 $xOy$ 中，抛物线 $\Gamma\colon y^2 = 2px$ $(p > 0)$
    的焦点为 $F$，过 $\Gamma$ 上一点 $P$（异于 $O$）作 $\Gamma$ 的切线，
    与 $y$ 轴交于点 $Q$. 若 $|FP| = 2$，$|FQ| = 1$，
    则向量 $\overrightarrow{OP}$ 与 $\overrightarrow{OQ}$ 的数量积为 \blankline
  \item 一颗质地均匀的正方体骰子，六个面上分别标有点数 $1, 2, 3, 4, 5, 6$.
    随机地投掷该骰子三次（各次抛掷结果相互独立），所得的点数依次为
    $a_1$, $a_2$, $a_3$，则事件
    “$|a_1 - a_2| + |a_2 - a_3| + |a_3 - a_1| = 6$”发生的概率为 \blankline
  \item 设有理数 $r = \dfrac{p}{q} \in (0,1)$，其中 $p$, $q$ 为互素的正整数，
    且 $pq$ 整除 $3600$. 这样的有理数 $r$ 的个数为 \blankline
\end{enumerate}

\section{解答题：本大题共 3 个小题，满分 56 分. 解答应写出文字说明、证明过程或演算步骤.}

\begin{enumerate}[resume]
  \item (本题满分 16 分) 已知复数列 $\{z_n\}$ 满足：
    $z_1 = \dfrac{\sqrt{3}}{2}$，$z_{n+1} = \bar{z}_n(1 + z_n\ii)$
    $(n = 1, 2, \cdots)$，其中 $\ii$ 为虚数单位. 求 $z_{2021}$ 的值.
  \item (本题满分 20 分) 在平面直角坐标系中，函数
    $y = \dfrac{x+1}{|x|+1}$ 的图象上有三个不同的点位于直线 $l$ 上，
    且这三点的横坐标之和为 $0$. 求 $l$ 的斜率的取值范围.
  \item (本题满分 20 分) 如图，正方体 $ABCD\text{-}EFGH$ 的棱长为 $2$，
    在正方形 $ABFE$ 的内切圆上任取一点 $P_1$，在正方形 $BCGF$
    的内切圆上任取一点 $P_2$，在正方形 $EFGH$ 的内切圆上任取一点 $P_3$.
    求 $|P_1P_2| + |P_2P_3| + |P_3P_1|$ 的最小值与最大值.
\end{enumerate}

\end{document}